\subsection*{CU-11: Borrar la cuenta}
\addcontentsline{toc}{subsection}{CU-11: Borrar la cuenta}
\label{cu-11}

\begin{fichacu}{CU-11}{Borrar la cuenta}
  \crudFila{Responsable}{José Eduardo Prior Hernández}
  \crudFila{Actor principal}{Jugador}
  \crudFila{Actores de sistema}{Servidor de partidas, Servicio de correo, Base de datos}
  \crudFila{Prototipos}{9f, 9g}
  \crudFila{Objetivo}{Permitir al Jugador dar de baja su cuenta con una confirmación en dos pasos, conservando durante catorce días la posibilidad de recuperarla y sin romper el historial de las partidas que compartió con otros jugadores.}
  \crudFila{Disparador}{El Jugador selecciona ``borrar cuenta'' en la \texttt{GUI\_\allowbreak{}AccountSettings}.}
  \textbf{Precondiciones} & \cuCelda{\begin{cureglas}
      \item \textbf{\mbox{PRE-01}} El Jugador tiene una \textit{Sesion} abierta y su token es válido.
      \item \textbf{\mbox{PRE-02}} El \textit{Usuario} tiene \texttt{tipo\_\allowbreak{}cuenta} en \texttt{REGISTRADA} y \texttt{estado\_\allowbreak{}cuenta} en \texttt{ACTIVA}.
      \item \textbf{\mbox{PRE-03}} El Servidor de partidas está en ejecución y tiene conexión disponible con la Base de datos.
    \end{cureglas}} \\ \hline
  \textbf{Flujo normal} & \cuCelda{\begin{cuenum}[start=1]
      \item El Jugador selecciona ``borrar cuenta'' en la \texttt{GUI\_\allowbreak{}AccountSettings}.
      \item El SJM despliega el primer paso de la \texttt{GUI\_\allowbreak{}DeleteAccount}: una lista de lo que se pierde —elo, monedas, objetos, amigos y posición en el ranking—, el aviso de que la cuenta se puede recuperar durante catorce días y un campo para la contraseña actual, con las opciones ``Continuar'' y ``Cancelar''. \textit{(Ver \mbox{FA-01})}
      \item El Jugador captura su contraseña y da click en ``Continuar''.
      \item El SJM despliega el segundo paso, que pide escribir la palabra de confirmación indicada en pantalla, con las opciones ``Borrar mi cuenta'' y ``Cancelar'' (\mbox{RN-01}).
      \item El Jugador escribe la palabra y da click en ``Borrar mi cuenta''.
      \item El SJM verifica que la palabra coincida exactamente. \textit{(Ver \mbox{FA-02})}
      \item El SJM envía el mensaje \texttt{DELETE\_\allowbreak{}ACCOUNT\_\allowbreak{}REQUEST} con la contraseña y el token de sesión. \textit{(Ver \mbox{EX-03}, \mbox{EX-04}, \mbox{EX-05})}
      \item El Servidor de partidas verifica que el token corresponda a una \textit{Sesion} abierta. \textit{(Ver \mbox{FA-03})}
    \end{cuenum}} \\
   & \cuCelda{\begin{cuenum}[start=9]
      \item El Servidor de partidas deriva el resumen de la contraseña recibida y lo compara contra \texttt{contrasena\_\allowbreak{}hash}. \textit{(Ver \mbox{FA-04})}
      \item El Servidor de partidas verifica que el \textit{Usuario} no tenga ninguna \textit{Participacion} sin terminar en una partida en línea (\mbox{RN-04}). \textit{(Ver \mbox{FA-05})}
      \item El Servidor de partidas verifica que el \texttt{rol} del \textit{Usuario} no sea \texttt{ADMINISTRADOR} (\mbox{RN-06}). \textit{(Ver \mbox{FA-06})}
      \item El Servidor de partidas inicia una transacción sobre la Base de datos.
      \item El Servidor de partidas cambia el \texttt{estado\_\allowbreak{}cuenta} del \textit{Usuario} a \texttt{ELIMINADA} y escribe \texttt{fecha\_\allowbreak{}baja} con la fecha y hora actuales (\mbox{RN-02}). \textit{(Ver \mbox{EX-01})}
      \item El Servidor de partidas retira al \textit{Usuario} de la \textit{ColaEmparejamiento} y de cualquier \textit{SalaParticipante}, si estuviera.
      \item El Servidor de partidas elimina las \textit{Solicitud} de amistad pendientes, enviadas y recibidas, y marca como expiradas las \textit{Invitacion} pendientes del \textit{Usuario}.
      \item Si el \texttt{rol} del \textit{Usuario} es \texttt{MODERADOR}, el Servidor de partidas devuelve a \texttt{PENDIENTE} los \textit{Reporte} que tuviera \texttt{EN\_\allowbreak{}REVISION}.
    \end{cuenum}} \\
   & \cuCelda{\begin{cuenum}[start=17]
      \item El Servidor de partidas cierra todas las \textit{Sesion} abiertas del \textit{Usuario}, con \texttt{motivo\_\allowbreak{}cierre} igual a \texttt{BAJA\_\allowbreak{}CUENTA}.
      \item El Servidor de partidas confirma la transacción. \textit{(Ver \mbox{EX-02})}
      \item El Servidor de partidas registra la baja en la \textit{BitacoraAcceso} con el resultado \texttt{CUENTA\_\allowbreak{}ELIMINADA}.
      \item El Servidor de partidas solicita al Servicio de correo un aviso al \texttt{correo} del \textit{Usuario} con la fecha límite de recuperación, fuera de la transacción. \textit{(Ver \mbox{EX-06})}
      \item El Servidor de partidas responde con \texttt{DELETE\_\allowbreak{}ACCOUNT\_\allowbreak{}OK} y la fecha límite de recuperación.
      \item El SJM descarta el token y todos los datos del perfil guardados en memoria, y despliega la \texttt{GUI\_\allowbreak{}Login} con un aviso que indica hasta qué fecha puede recuperar la cuenta iniciando sesión.
      \item Fin del caso de uso.
    \end{cuenum}} \\ \hline
  \textbf{Flujos alternos} & \cuCelda{\textbf{\mbox{FA-01}: El Jugador cancela en cualquiera de los dos pasos}\par
    \begin{cuenum}[start=1]
      \item El Jugador selecciona ``Cancelar''.
      \item El SJM descarta lo capturado sin enviarlo y vuelve a la \texttt{GUI\_\allowbreak{}AccountSettings}.
      \item Fin del caso de uso.
    \end{cuenum}} \\
   & \cuCelda{\textbf{\mbox{FA-02}: La palabra de confirmación no coincide}\par
    \begin{cuenum}[start=1]
      \item El SJM mantiene deshabilitado ``Borrar mi cuenta'' mientras la palabra no coincida y lo indica junto al campo.
      \item El flujo continúa en el paso 5 del FN.
    \end{cuenum}} \\
   & \cuCelda{\textbf{\mbox{FA-03}: La sesión ya no es válida}\par
    \begin{cuenum}[start=1]
      \item El Servidor de partidas responde con \texttt{SESSION\_\allowbreak{}INVALID}.
      \item El SJM muestra la \texttt{GUI\_\allowbreak{}MessageWarning} indicando que la sesión terminó y despliega la \texttt{GUI\_\allowbreak{}Login}. La cuenta no se borró.
      \item Fin del caso de uso.
    \end{cuenum}} \\
   & \cuCelda{\textbf{\mbox{FA-04}: La contraseña es incorrecta}\par
    \begin{cuenum}[start=1]
      \item El Servidor de partidas incrementa \texttt{intentos\_\allowbreak{}fallidos} y aplica el bloqueo escalonado del \mbox{RN-03} del \mbox{CU-01}, igual que en el \mbox{FA-05} del \mbox{CU-09}.
      \item El Servidor de partidas registra el intento en la \textit{BitacoraAcceso} con el resultado \texttt{CONTRASENA\_\allowbreak{}INCORRECTA}.
      \item El Servidor de partidas responde con \texttt{DELETE\_\allowbreak{}ACCOUNT\_\allowbreak{}FAILED} y el motivo \texttt{CONTRASENA\_\allowbreak{}INCORRECTA}.
      \item El SJM vuelve al primer paso con el campo de contraseña vacío.
      \item El flujo continúa en el paso 3 del FN.
    \end{cuenum}} \\
   & \cuCelda{\textbf{\mbox{FA-05}: El Jugador tiene una partida en curso}\par
    \begin{cuenum}[start=1]
      \item El Servidor de partidas responde con \texttt{DELETE\_\allowbreak{}ACCOUNT\_\allowbreak{}FAILED} y el motivo \texttt{PARTIDA\_\allowbreak{}EN\_\allowbreak{}CURSO}.
      \item El SJM muestra la \texttt{GUI\_\allowbreak{}MessageWarning} indicando que primero debe terminar o abandonar la partida, con las opciones ``Volver a la partida'' y ``Aceptar''.
      \item Si el Jugador selecciona ``Volver a la partida'', el flujo continúa en el \mbox{CU-26}. En caso contrario, vuelve a la \texttt{GUI\_\allowbreak{}AccountSettings}.
      \item Fin del caso de uso.
    \end{cuenum}} \\
   & \cuCelda{\textbf{\mbox{FA-06}: La cuenta es de administración}\par
    \begin{cuenum}[start=1]
      \item El Servidor de partidas responde con \texttt{DELETE\_\allowbreak{}ACCOUNT\_\allowbreak{}FAILED} y el motivo \texttt{CUENTA\_\allowbreak{}DE\_\allowbreak{}ADMINISTRACION}.
      \item El SJM indica que una cuenta de administración no puede darse de baja desde el juego.
      \item Fin del caso de uso.
    \end{cuenum}} \\
   & \cuCelda{\textbf{\mbox{FA-07}: El Jugador es un invitado}\par
    \begin{cuenum}[start=1]
      \item El SJM no muestra la opción ``borrar cuenta'' a una cuenta de invitado, que no tiene contraseña con la que confirmar.
      \item En su lugar, el SJM ofrece ``olvidar este invitado'', que descarta el token guardado en el dispositivo sin tocar la base de datos: la cuenta de invitado caducará por inactividad conforme al \mbox{RN-10} del \mbox{CU-06}.
      \item Fin del caso de uso.
    \end{cuenum}} \\
   & \cuCelda{\textbf{\mbox{FA-08}: Transcurren los catorce días sin recuperar la cuenta}\par
    \begin{cuenum}[start=1]
      \item El proceso programado del Servidor de partidas encuentra el \textit{Usuario} con \texttt{estado\_\allowbreak{}cuenta} en \texttt{ELIMINADA} y \texttt{fecha\_\allowbreak{}baja} anterior a catorce días.
      \item El Servidor de partidas sustituye \texttt{nickname} y \texttt{correo} por valores anónimos únicos, deja \texttt{contrasena\_\allowbreak{}hash} y \texttt{contrasena\_\allowbreak{}sal} en nulo, elimina sus \textit{EnlaceRed}, invalida su \textit{Avatar} y elimina sus \textit{Amistad}, \textit{Silencio} y \textit{Bloqueo} (\mbox{RN-03}).
      \item El Servidor de partidas conserva el \textit{Usuario} anonimizado y todo lo que otros jugadores comparten con él: \textit{Participacion}, \textit{Jugada}, \textit{Mensaje}, \textit{Reporte} y \textit{Sancion}.
      \item El Servidor de partidas registra la anonimización en la \textit{BitacoraAcceso} con el resultado \texttt{CUENTA\_\allowbreak{}ANONIMIZADA}.
      \item Fin del caso de uso.
    \end{cuenum}} \\ \hline
  \textbf{Excepciones} & \cuCelda{\textbf{\mbox{EX-01}: Fallo al escribir en la base de datos}\par
    \begin{cuenum}[start=1]
      \item El Servidor de partidas revierte la transacción completa, de modo que no quede una cuenta eliminada con sesiones abiertas o con solicitudes pendientes.
      \item El Servidor de partidas registra la excepción en la bitácora del servidor con un identificador de incidencia y responde con \texttt{SERVER\_\allowbreak{}ERROR}.
      \item El SJM muestra la \texttt{GUI\_\allowbreak{}MessageError} con el identificador. La cuenta no se borró.
      \item Fin del caso de uso.
    \end{cuenum}} \\
   & \cuCelda{\textbf{\mbox{EX-02}: Fallo al confirmar la transacción}\par
    \begin{cuenum}[start=1]
      \item El Servidor de partidas trata la transacción como fallida y registra la excepción señalando que el estado de la baja es indeterminado.
      \item El Servidor de partidas responde con \texttt{SERVER\_\allowbreak{}ERROR}.
      \item El SJM muestra la \texttt{GUI\_\allowbreak{}MessageError} indicando que no fue posible confirmar la baja y que, si al volver a entrar se le ofrece recuperar la cuenta, la baja sí se aplicó.
      \item Fin del caso de uso.
    \end{cuenum}} \\
   & \cuCelda{\textbf{\mbox{EX-03}: Se agota el tiempo de espera de la respuesta}\par
    \begin{cuenum}[start=1]
      \item El SJM detecta que transcurrió el tiempo límite sin recibir un mensaje completo.
      \item El SJM no reenvía la solicitud, descarta el token local y despliega la \texttt{GUI\_\allowbreak{}Login}, indicando que la baja pudo haberse aplicado y que lo sabrá al volver a iniciar sesión.
      \item Fin del caso de uso.
    \end{cuenum}} \\
   & \cuCelda{\textbf{\mbox{EX-04}: La conexión se interrumpe durante el intercambio}\par
    \begin{cuenum}[start=1]
      \item El SJM detecta el cierre inesperado del socket antes de recibir respuesta.
      \item El flujo continúa en el paso 2 del \mbox{EX-03}.
    \end{cuenum}} \\
   & \cuCelda{\textbf{\mbox{EX-05}: El mensaje recibido está malformado}\par
    \begin{cuenum}[start=1]
      \item El SJM descarta el mensaje, cierra la conexión y registra en la bitácora local el contenido truncado.
      \item El flujo continúa en el paso 2 del \mbox{EX-03}.
    \end{cuenum}} \\
   & \cuCelda{\textbf{\mbox{EX-06}: El servicio de correo no está disponible}\par
    \begin{cuenum}[start=1]
      \item El Servidor de partidas agota los reintentos previstos hacia el Servicio de correo.
      \item El Servidor de partidas conserva la baja, que ya está confirmada, y registra el fallo del aviso en la bitácora del servidor.
      \item El flujo continúa en el paso 21 del FN. El aviso en pantalla del paso 22 sigue informando la fecha límite de recuperación.
    \end{cuenum}} \\ \hline
  \textbf{Postcondiciones} & \cuCelda{\begin{cureglas}
      \item \textbf{\mbox{POST-01}} Si el caso termina por el FN, el \textit{Usuario} tiene \texttt{estado\_\allowbreak{}cuenta} en \texttt{ELIMINADA} y \texttt{fecha\_\allowbreak{}baja} con la fecha de la baja.
      \item \textbf{\mbox{POST-02}} Si el caso termina por el FN, ninguna \textit{Sesion} de la cuenta sigue abierta y el Jugador no está en ninguna cola ni sala.
      \item \textbf{\mbox{POST-03}} Si el caso termina por el FN, no queda ninguna \textit{Solicitud} de amistad pendiente del \textit{Usuario}.
      \item \textbf{\mbox{POST-04}} Si el caso termina por el FN, la cuenta no aparece en el ranking, en las búsquedas ni en las listas de amigos de otros jugadores, pero sus partidas pasadas siguen visibles en el historial de sus rivales.
      \item \textbf{\mbox{POST-05}} Durante los catorce días siguientes, el \texttt{nickname} y el \texttt{correo} siguen ocupados y la cuenta puede recuperarse con el \mbox{FA-17} del \mbox{CU-01}.
      \item \textbf{\mbox{POST-06}} Si el caso termina por el \mbox{FA-08}, el \textit{Usuario} está anonimizado: su \texttt{nickname} y su \texttt{correo} quedan libres y ya no puede recuperarse.
      \item \textbf{\mbox{POST-07}} Ninguna fila de \textit{Usuario} se elimina físicamente en ningún flujo.
    \end{cureglas}} \\ \hline
  \textbf{Reglas de negocio} & \cuCelda{\begin{cureglas}
      \item \textbf{\mbox{RN-01}} Borrar la cuenta exige confirmación en dos pasos: la contraseña actual y la palabra de confirmación escrita. Es la confirmación más fuerte del sistema, junto con la del cambio de nombre.
      \item \textbf{\mbox{RN-02}} El borrado es lógico: la cuenta pasa a \texttt{ELIMINADA} y se conserva durante catorce días desde \texttt{fecha\_\allowbreak{}baja}, durante los cuales puede recuperarse iniciando sesión (\mbox{D-07}).
      \item \textbf{\mbox{RN-03}} Transcurridos los catorce días, la cuenta se anonimiza en lugar de eliminarse. Los datos personales desaparecen y los identificadores quedan libres, pero la fila se conserva porque otros registros dependen de ella.
      \item \textbf{\mbox{RN-04}} No puede borrarse una cuenta con una partida en línea sin terminar: la partida pertenece también a los demás jugadores y debe cerrarse primero.
      \item \textbf{\mbox{RN-05}} Las partidas, jugadas, mensajes, reportes y sanciones del \textit{Usuario} no se eliminan en ningún momento. Borrarlos alteraría el historial y las pruebas de otros jugadores.
      \item \textbf{\mbox{RN-06}} Una cuenta con rol de administrador no puede borrarse desde el juego.
      \item \textbf{\mbox{RN-07}} Las cuentas de invitado no usan este caso: no tienen contraseña, y su baja equivale a olvidar el token del dispositivo.
    \end{cureglas}} \\ \hline
  \textbf{Observación} & \cuCelda{\textit{Sobre la contradicción con el prototipo.}\par
    El prototipo 9f dice, bajo el botón de borrar cuenta, que ``no hay vuelta atrás'', mientras que \texttt{descripcion.md} promete catorce días de recuperación. Se resolvió a favor del documento de diseño (\mbox{D-07}), y el texto del prototipo debe corregirse para que diga la verdad: es precisamente el aviso que el jugador lee antes de decidir. La recuperación no necesita un caso de uso propio, porque ocurre al iniciar sesión dentro del plazo y se describe en el \mbox{FA-17} del \mbox{CU-01}.} \\ \hline
  \textbf{Observación} & \cuCelda{\textit{Sobre por qué se anonimiza en lugar de eliminar.}\par
    Eliminar físicamente la fila de \textit{Usuario} arrastraría por clave foránea sus participaciones, sus jugadas y sus mensajes, que son también el historial, las repeticiones y las pruebas de reportes de otros jugadores. Poner esas claves en nulo evitaría el arrastre, pero dejaría partidas con un participante inexistente. Anonimizar conserva la integridad referencial y cumple el propósito de la baja: que ya no quede nada que identifique a la persona.} \\ \hline
  \textbf{Datos que toca} & \cuCelda{\begin{cudatos}
      \item \textbf{\textit{Sesion}:} lee por token; cierra todas \textit{(pasos 8, 17)}
      \item \textbf{\textit{Usuario}:} lee \texttt{contrasena\_\allowbreak{}hash}, \texttt{contrasena\_\allowbreak{}sal}, \texttt{rol} \textit{(pasos 9, 11)}
      \item \textbf{\textit{Usuario}:} escribe \texttt{estado\_\allowbreak{}cuenta}, \texttt{fecha\_\allowbreak{}baja}; en el \mbox{FA-08} anonimiza \texttt{nickname}, \texttt{correo}, \texttt{contrasena\_\allowbreak{}hash}, \texttt{contrasena\_\allowbreak{}sal} \textit{(pasos 13, \mbox{FA-08})}
      \item \textbf{\textit{Participacion}:} lee, para comprobar que no haya partida en curso \textit{(paso 10)}
      \item \textbf{\textit{ColaEmparejamiento}, \textit{SalaParticipante}:} elimina la del usuario \textit{(paso 14)}
      \item \textbf{\textit{Solicitud}:} elimina las pendientes \textit{(paso 15)}
      \item \textbf{\textit{Invitacion}:} marca como expiradas las pendientes \textit{(paso 15)}
      \item \textbf{\textit{Reporte}:} devuelve a \texttt{PENDIENTE} los que tenía en revisión \textit{(paso 16)}
    \end{cudatos}} \\
   & \cuCelda{\begin{cudatos}
      \item \textbf{\textit{EnlaceRed}, \textit{Amistad}, \textit{Silencio}, \textit{Bloqueo}:} elimina, en el \mbox{FA-08} \textit{(paso \mbox{FA-08})}
      \item \textbf{\textit{Avatar}:} invalida, en el \mbox{FA-08} \textit{(paso \mbox{FA-08})}
      \item \textbf{\textit{BitacoraAcceso}:} inserta \textit{(pasos 19, \mbox{FA-08})}
    \end{cudatos}} \\ \hline
\end{fichacu}
\clearpage
